% GENERATED by tools/build_m05_code_sheet.py -- do not edit.
\begin{codebox}{Box 6}{The US airports: Leiden, and a map}
import seaborn as sns
import pandas as pd
members = g_air.community_leiden("modularity", n_iterations=-1).membership
top = pd.Series(members).value_counts().index[:4]
sns.scatterplot(x=lon, y=lat, color="#bdbdbd", s=14)
sns.scatterplot(
    x=lon, y=lat, hue=members, hue_order=top, palette=palette[:4], s=14
)
plt.show()
\end{codebox}
\begin{linenotes}
\item[1--2] \texttt{seaborn} draws statistical pictures, and \texttt{pandas} works with tables. \texttt{sns} and \texttt{pd} are their usual short names.
\item[3] Box 2, line 1, again, on the US airport network \texttt{g\_air} (540 airports, two linked when a flight connects them), which is ready-made. We keep the group of each airport as \texttt{members}.
\item[4] \texttt{pd.Series(members).value\_counts()} counts how many airports each group has, biggest group first. \texttt{.index[:4]} keeps the numbers of the four biggest groups. We call them \texttt{top}.
\item[5] \texttt{lon} and \texttt{lat} are the longitude and the latitude of each airport, also ready-made. \texttt{sns.\allowbreak scatterplot} puts one dot per airport at that place: a rough map of the US. This first call paints all of them light grey; \texttt{s=14} is the dot size.
\item[6--8] The second call paints over the grey, but only the four groups in \texttt{top}: \texttt{hue=members} says that the colour follows the group, \texttt{hue\_order=top} limits it to those four, and \texttt{palette[:4]} is the first four colours of the palette. Every other airport stays light grey: it is ``the rest''.
\item[9] Show the map.
\end{linenotes}
